\section{模型假设}

题目给出的示向误差、接收距离、移动速度和操作耗时直接作为模型参数，
不再重复列为假设。在此基础上，仅作以下简化：

\begin{enumerate}
  \item 将目标区域视为平面欧氏空间，机器狗的位置坐标和当前朝向没有额外测量误差。这样可以直接用点间距离计算移动路程，并用示向度误差范围约束源位置；
  \item 机器狗在相邻停驻点之间沿线段匀速运动，转向和通信耗时忽略不计。题目只规定了移动速度和各项操作时间，故总耗时由移动时间与操作时间相加计算；
  \item 一次测试期间，干扰源的位置、频道、有效接收半径和发射方向保持不变。该假设使不同时刻取得的观测能够共同约束同一个源位置；
  \item 接口反馈能够正确区分示向度、近距离信号和无信号三类结果。模型只处理题目给定的示向误差以及接口保留两位小数造成的舍入误差，不另设漏检或误报概率。
\end{enumerate}

本文不把同一位置的重复读数当作独立随机样本，也不通过多次平均缩小
题设误差界。问题一采用 \(1\degree\) 的误差半角；问题二至四考虑两位
小数的舍入误差，取 \(1.0051\degree\) 作为计算半角。
